\documentclass[11pt]{article}
\usepackage[T1]{fontenc}
\usepackage[utf8]{inputenc}
\usepackage{lmodern,amsmath,amssymb,amsthm,mathtools}
\usepackage[margin=28mm]{geometry}
\usepackage{microtype,xcolor}
\usepackage[authoryear,round]{natbib}
\usepackage{xurl}
\usepackage[colorlinks=true,allcolors=blue!50!black]{hyperref}
\newtheorem{theorem}{Theorem}[section]
\newtheorem{proposition}[theorem]{Proposition}
\newtheorem{lemma}[theorem]{Lemma}
\newtheorem{corollary}[theorem]{Corollary}
\theoremstyle{definition}\newtheorem{definition}[theorem]{Definition}
\theoremstyle{remark}\newtheorem{remark}[theorem]{Remark}
\newcommand{\OO}{\mathcal O}
\newcommand{\PP}{\mathbf P}
\newcommand{\FF}{\mathbf F}
\newcommand{\Pic}{\operatorname{Pic}}
\newcommand{\NS}{\operatorname{NS}}
\newcommand{\End}{\operatorname{End}}
\newcommand{\Hom}{\operatorname{Hom}}
\newcommand{\rk}{\operatorname{rank}}
\newcommand{\CH}{\operatorname{CH}}
\newcommand{\SB}{\operatorname{SB}}
\newcommand{\st}{\mathrm{st}}
\newcommand{\workbenchurl}{https://github.com/KokunoYumeto/schemes-and-moduli-workbench}
\newcommand{\recordurl}{https://github.com/KokunoYumeto/mathematics-workbenches}
\title{Grothendieck's generic Frobenius ranks for primitive odd polarizations, after Cartier, Sun, and Karpenko}
\author{}
\date{23 November 2026}
\begin{document}
\maketitle
\begin{abstract}
Grothendieck's crystalline-cohomology notes ask how the generic stable Frobenius rank of hypersurface sections depends on their degree. We determine the relative rank for every fixed nonordinary abelian surface in characteristic two with a globally generated polarization of type $(1,\ell)$, for $\ell=3,5,7,11,13,17$. A general smooth curve has genus $\ell+1$ and $2$-rank $f_A+\ell-1$, where $f_A$ is the $2$-rank of the surface. For every odd prime the same conclusion follows from one explicit maximal-rank condition on the dual Cartier pencil. Cartier multiplication and Sun's Frobenius filtration give a nonzero theta-equivariant tensor. Karpenko's incompressibility argument forces its radical map to be dominant; an incidence calculation makes it birational. Exact polynomial certificates establish the required dual rank for the six listed primes. The type $(1,3)$ example has a Fermat cubic as its nonbijective locus. Global generation remains an explicit hypothesis, and no monodromy conclusion is inferred. The mathematics and text were produced and checked by GPT-6 Astra Pro with the scopes stated at the end; no human has checked the mathematics.
\end{abstract}

\section{The question and the rank theorem}

This paper determines the surviving Frobenius contribution of curves in primitive prime-degree linear systems on a fixed abelian surface. In \S8.5.8 of his crystalline notes, Grothendieck writes: ``Nous d\'esirons montrer que pour $d$ assez grand, ce nombre est non nul'' \citep[cote~4, archive sheets~7--8, PDF images~8--9]{Grothendieck}. The number is the generic count of $p$-adic unit Frobenius eigenvalues in evanescent cohomology. His immediate goal is eventual positivity; here we calculate the relative stable rank exactly for specified primitive systems. The result here concerns complete linear systems of primitive type. It does not require identifying such a system with a multiple of a principal polarization. That distinction matters: their section spaces and theta representations have different dimensions.

Let $k$ be algebraically closed of characteristic two. If $A$ is an abelian surface, its $2$-rank $f_A$ is defined by $|A[2](k)|=2^{f_A}$. The surface is \emph{nonordinary} when $f_A<2$. For a smooth connected projective curve $C$, let $f_C$ denote the $2$-rank of its Jacobian. An ample line bundle $L$ has type $(1,\ell)$, for an odd prime $\ell$, when the kernel of
\[
 \lambda_L:A\longrightarrow A^\vee,\qquad
 x\longmapsto t_x^*L\otimes L^{-1},
\]
is isomorphic to $(\mathbb Z/\ell)^2$. Equivalently, its prime-to-characteristic polarization type is $(1,\ell)$. Then
\[
 N:=h^0(A,L)=\chi(L)=\frac{L^2}{2}=\ell.
\]
Global generation means that $H^0(A,L)\otimes\OO_A\to L$ is surjective at every point.

For a smooth member $C\in|L|$, restriction on $H^1(\OO)$ is injective. Define
\[
 Q_C=\operatorname{coker}\bigl(H^1(A,\OO_A)\longrightarrow
 H^1(C,\OO_C)\bigr),\qquad
 r(C)=\dim_k(Q_C)_{\st},
\]
where $V_{\st}=\bigcap_{j\ge0}\operatorname{im}F^j$ for a finite-dimensional space with semilinear Frobenius. We will prove
\begin{equation}\label{eq:dimensions}
 g(C)=\ell+1,\qquad \dim Q_C=\ell-1,\qquad r(C)=f_C-f_A.
\end{equation}

\begin{theorem}[Primitive prime polarizations]\label{thm:main}
Let $A$ be a fixed nonordinary abelian surface over $k$, and let $L$ be ample and globally generated of type $(1,\ell)$, where
\[
 \ell\in\{3,5,7,11,13,17\}.
\]
There is a nonempty open of smooth connected curves $C\in|L|$ on which
\begin{equation}\label{eq:main}
 g(C)=\ell+1,\qquad f_C=f_A+\ell-1,\qquad r(C)=\ell-1.
\end{equation}
Thus Frobenius on $Q_C$ is bijective. Over a perfect field, the same conclusion holds for geometric generic rank after algebraic closure.
\end{theorem}

In words, a general curve in each of these systems has exactly the defect $2-f_A$ already present in the surface. Its relative cohomology contributes the largest possible stable rank. For instance, type $(1,3)$ gives genus four and $2$-rank two on a surface of $2$-rank zero, or $2$-rank three on a surface of $2$-rank one. This is a fixed-surface answer to the stable-rank part of Grothendieck's programme, rather than a statement about a general abelian surface in a moduli space.

Here is the exact condition that extends the theorem to an arbitrary odd prime. Let $F_A:A\to B=A^{(2)}$ be relative Frobenius, put $H=L^{(2)}$ and $W=H^0(A,L)$, and choose an invariant volume form. Relative Cartier trace of products defines
\begin{equation}\label{eq:beta-intro}
 \beta:\bigwedge^2W\longrightarrow H^0(B,H)=W^{(2)}.
\end{equation}
The superscript $(2)$ denotes the Frobenius twist, not a tensor square. The equality $F_A^*H\simeq L^2$ explains the target. For $y\in(W^{(2)})^*$, write $B_y(v,w)=y(\beta(v,w))$.

\begin{theorem}[The conditional rank criterion]\label{thm:conditional}
Retain the geometric hypotheses of Theorem~\ref{thm:main}, but allow any odd prime $\ell$. If the alternating form $B_y$ has rank $\ell-1$ for generic $y$, then all conclusions of \eqref{eq:main} hold. The map \eqref{eq:beta-intro} is surjective without this additional rank hypothesis.
\end{theorem}

The additional hypothesis is explicit. In a theta basis $e_i$, indexed by $\mathbb Z/\ell$, there are coefficients
\[
 c_0=0,\qquad c_a=c_{-a},\qquad
 B_y(i,j)=c_{i-j}y_{(i+j)/2}.
\]
Division by two is in $\mathbb Z/\ell$. The coefficients are not all zero. We prove that every nonzero coefficient vector has the required generic rank for the six primes in Theorem~\ref{thm:main}; the general-prime assertion is left open.

\paragraph{Credit and context.}
The local Cartier formula and its multiplicativity are classical \citep[0, (2.1.4)--(2.1.13), pp.~514--515]{Illusie}; the relation between stable coherent cohomology and $p$-rank is part of \citet[Propositions~3.18 and 3.22--3.24, pp.~89, 99--101]{Oda}. The Frobenius filtration used below is the surface specialization of \citet[Definition~3.1, Lemmas~3.2--3.5 and Theorem~3.6, pp.~5--9]{Sun}. We prove the resulting slope-stability assertion directly, including its Euler-characteristic obstruction. Karpenko's prime-degree incompressibility theorem is known; the diagonal-trace specialization of his method is proved in full in Section~\ref{sec:dominance}, following \citet[Lemmas~2.21, 2.23 and Corollaries~2.22, 2.24, pp.~9--10]{Karpenko}. Odd theta descent and theta representation theory are used at the exact primary locations given below. The generic-rank conclusion combines these inputs with the two contraction matrices and universal coefficient certificates.

For type $(1,3)$ over $\mathbb C$, \citet[Introduction, Theorems~0.1--0.2, pp.~1--2]{CataneseSernesi} study the discriminant of singular divisors using Heisenberg symmetry and the Hesse pencil. The cubic in Proposition~\ref{prop:cubic} instead records nonbijectivity of relative Frobenius in characteristic two.

Recent AI-assisted work of \citet[Introduction, Theorems~A--C, and acknowledgements]{Zou} studies quasi-$F$-splitting and crystalline slopes of primitive symplectic varieties. Its Frobenius questions are neighboring context, not an input to the fixed linear-system argument here. The preprint of \citet[Introduction and Theorem~1.1, pp.~2--5]{OpenAI} uses characteristic-two matrix geometry in a proof about complex border determinantal complexity; it also does not address the primitive Cartier-rank question. The present finite certificates concern symbolic coefficient polynomials; they do not infer a theorem from random values or finitely many fields. No priority claim is made.

\paragraph{Scope and mechanism.}
Global generation is used in the divisor bound and in smoothing a nodal member. Oddness makes the theta kernel \'etale and its central-weight representations rigid. Nonordinarity makes the rank-four Cartier pairing alternating. Primal multiplication and the dual alternating pencil are different matrices; surjectivity of their total tensor does not by itself give either matrix maximal rank. The proof first establishes the dual rank, then uses a twist by a prime-degree division algebra to force dominance of the radical map. Linear fibres make that map birational. Finally, the derivative of a degree-$\ell$ polynomial proves that the induced Frobenius on the Euler bundle is generically invertible. No assertion about a common rank-drop wall, a Lang cover, or a full monodromy group is needed or obtained. The parameter space is the complete linear system; ambient hypersurface restrictions require their own section-surjectivity condition.

\section{Cohomology and the universal Euler bundle}\label{sec:euler}

We identify the relative cohomology before calculating its Frobenius. We use the standard vanishing and Riemann--Roch formulas for ample lines on abelian surfaces \citep[Theorems~9.11 and 9.14, pp.~131--133]{EGM}. In particular, on our surface an ample $L$ has $H^i(A,L)=0$ for $i>0$, and Serre duality gives $H^1(A,L^{-1})=0$.

\begin{lemma}[Stable parts and a hypersurface]\label{lem:cohomology}
For a smooth ample curve $C\in|L|$, the sequence
\begin{equation}\label{eq:curve-sequence}
 0\longrightarrow H^1(A,\OO_A)\longrightarrow H^1(C,\OO_C)
 \longrightarrow Q_C\longrightarrow0
\end{equation}
is Frobenius compatible, $\dim Q_C=N-1$, and $f_C=f_A+\dim(Q_C)_{\st}$. Every effective member of $|L|$ is connected and has arithmetic genus $N+1$.
\end{lemma}
\begin{proof}
The divisor sequence and $H^1(L^{-1})=0$ give the injection and identify $Q_C$ with the kernel of $H^2(L^{-1})\to H^2(\OO_A)$. Under Serre duality this last map is dual to $k\to H^0(L)$ sending $1$ to the nonzero defining section. It is surjective, so its kernel has dimension $N-1$. The same sequence in degrees zero and one gives $H^0(C,\OO_C)=k$ for every effective member. Adjunction, with $\omega_A\simeq\OO_A$, gives arithmetic genus $1+L^2/2=N+1$.

For completeness, stable parts are exact for semilinear endomorphisms over a perfect field. For sufficiently large $m$, the stable image and the kernel of $F^m$ give the Fitting decomposition into bijective and nilpotent summands. It is functorial: an intertwining map sends the nilpotent summand to the nilpotent summand, and a map from a bijective summand to a nilpotent one is zero. A short exact sequence therefore decomposes into two short exact sequences. The stable coherent dimension of a proper smooth curve or abelian variety equals its $2$-rank, so applying exactness to \eqref{eq:curve-sequence} proves the formula. One can also see this identification from Artin--Schreier: $H^1_{\mathrm{et}}(-,\FF_2)$ is the fixed space of $F$ on $H^1(\OO)$, and the dimension of this fixed space over $\FF_2$ is the stable dimension over $k$. On a bijective $d$-dimensional summand the additive map $v\mapsto F(v)-v$ is finite \'etale of degree $2^d$ and surjective; its kernel has $2^d$ points. On a nilpotent summand it is invertible by a finite geometric sum.
\end{proof}

\begin{lemma}[Universal Frobenius]\label{lem:euler}
Assume $A$ is nonordinary. Let $X=\PP(W)$ parametrize one-dimensional subspaces of $W$, and let $q:\mathcal C\to X$ be the universal divisor. The quotient of $R^1q_*\OO_{\mathcal C}$ by $H^1(A,\OO_A)\otimes\OO_X$ is
\[
 \mathcal E=\ker\bigl(W^*\otimes\OO_X(-1)\longrightarrow\OO_X\bigr)
 =\Omega_X^1.
\]
It carries a regular linearized Frobenius $\Phi:F_X^*\mathcal E\to\mathcal E$, where $F_X$ is absolute Frobenius. On a coefficient chart, with $s\in W$ representing its point, put $M_s(v)=\beta(s,v)$. Then
\begin{equation}\label{eq:euler-map}
 \Phi_s=M_s^t\big|_{(s^{[2]})^\perp}:
 (s^{[2]})^\perp\longrightarrow s^\perp.
\end{equation}
If $\rk M_s=N-1$ and $\ker M_s^t=ky_s$, then
\begin{equation}\label{eq:criterion}
 \Phi_s\text{ is invertible}\quad\Longleftrightarrow\quad
 \sum_i y_{s,i}s_i^2\ne0.
\end{equation}
The bundle and formula are defined even when the divisor is singular.
\end{lemma}
\begin{proof}
On $A\times X$ the universal ideal is $L^{-1}\boxtimes\OO_X(-1)$. Push its divisor sequence to $X$. The vanishing used above gives an exact sequence
\[
 0\to H^1(A,\OO_A)\otimes\OO_X\to R^1q_*\OO_{\mathcal C}
 \to H^2(A,L^{-1})\otimes\OO_X(-1)\to\OO_X\to0.
\]
The last map is evaluation after Serre duality. It is surjective at every point of $X$, since $s\ne0$, so its kernel is the Euler bundle as claimed. This also proves local freeness and base change for the displayed terms.

Frobenius preserves the ambient summand. Its formula on the ideal cohomology can be read directly from the divisor sequence: locally the ideal embeds by multiplication with $s$, and the square of $sh$ is $s\cdot(sh^2)$. Hence the induced map on ideal cohomology is $h\mapsto s h^2$. Serre duality identifies its transpose with relative Cartier trace of $sv$, namely $M_s(v)$. On the parameter space the Frobenius pullback of the evaluation row $s$ is the row $s^{[2]}$. These observations give \eqref{eq:euler-map}, including the indicated source rather than $s^\perp$ again. Alternation yields $M_ss=0$, so the image lies in $s^\perp$. When the kernel of $M_s^t$ is the stated line, restriction to a hyperplane is injective precisely when that line is not in the hyperplane. This is \eqref{eq:criterion}; the source and target both have dimension $N-1$.
\end{proof}

\section{Odd Frobenius pushforwards and Cartier multiplication}\label{sec:cartier}

We now prove that the geometric Cartier tensor cannot vanish, using slope stability and a divisor bound. For a torsion-free sheaf $E$ and an ample line $P$ on a surface, its slope is $\mu_P(E)=c_1(E)\cdot P/\rk E$. Stability means that every proper nonzero subsheaf of smaller rank has strictly smaller slope.

\begin{lemma}[Odd Euler characteristic and slope stability]\label{lem:stable}
For any line bundle $L$ on an abelian surface with odd $\chi(L)$, the rank-four bundle $F_{A*}L$ is slope stable with respect to every ample polarization on $B=A^{(2)}$.
\end{lemma}
\begin{proof}
The canonical Frobenius filtration has graded terms $L$, $L\otimes\Omega_A^1$, and $L\otimes\bigwedge^2\Omega_A^1$. This is the case of dimension and characteristic two of Sun's filtration at the cited locations. In this case it also has a direct construction. On $A\times_B A$ the diagonal ideal has local generators $\epsilon_1,\epsilon_2$ with $\epsilon_1^2=\epsilon_2^2=0$. Its successive degrees have ranks $1,2,1$ and transform as $\OO_A,\Omega_A^1,\bigwedge^2\Omega_A^1$. Tensoring with the line on the second factor gives the asserted filtration. Since the cotangent bundle of an abelian surface is trivial,
\[
 U=F_A^*F_{A*}L\otimes L^{-1}
\]
has a filtration by four copies of $\OO_A$. It is semistable of slope zero for any polarization. Consequently $F_{A*}L$ is semistable: a subsheaf of larger slope would pull back to one of larger slope in $U$.

Suppose a saturated proper subsheaf $Q\subset F_{A*}L$ has equal slope, and write $q=\rk Q\in\{1,2,3\}$. Put $S=F_A^*Q\otimes L^{-1}$ and $T=U/S$. Flatness of relative Frobenius makes $T$ torsion free. Both $S$ and $T$ have slope zero. Intersect $S$ with the trivial-quotient filtration of $U$. Every nonzero successive quotient is a rank-one subsheaf $I_Z(-D)\subset\OO_A$, where $D$ is effective and $Z$ is zero dimensional. Their degrees are nonpositive and sum to zero. Thus every $D$ vanishes. It follows that $c_1(S)=0$ and $\chi(S)=-\sum\operatorname{length}Z\le0$.

The injection $T^*\hookrightarrow U^*$ gives the same inequality $\chi(T^*)\le0$. On a smooth surface $T^{**}$ is locally free, and Serre duality with trivial canonical line gives $\chi(T^{**})=\chi(T^*)$. Therefore
\[
 \chi(T)=\chi(T^*)-\operatorname{length}(T^{**}/T)\le0.
\]
As $\chi(U)=4\chi(\OO_A)=0$, additivity forces $\chi(S)=0$. Let $l=c_1(L)$. The determinant equality gives $F_A^*c_1(Q)=q l$. Riemann--Roch, $\deg F_A=4$, and the formula for tensoring by $L^{-1}$ yield
\[
 0=\chi(S)=4\chi(Q)-F_A^*c_1(Q)\cdot l+q\frac{l^2}{2}
 =4\chi(Q)-q\chi(L).
\]
Oddness of $\chi(L)$ forces $4\mid q$, a contradiction. This proves strict stability, including the possible zero-dimensional defects of $Q$ and its quotient.
\end{proof}

\begin{lemma}[Divisor bound]\label{lem:divisor}
Let $H$ be ample and globally generated on an abelian surface, with odd $N=h^0(H)>1$. If $h^0(D)\ge N$ for a line bundle $D$, then $D\cdot H\ge2N$.
\end{lemma}
\begin{proof}
Choose an effective representative of $D$. Every effective divisor on an abelian variety is nef: a suitable translate contains no chosen curve, so its intersection with that curve is nonnegative. If $D^2>0$, Hodge index excludes a curve orthogonal to $D$, since such a curve would have negative square. Thus $D$ is ample. Vanishing and Riemann--Roch give $D^2=2h^0(D)\ge2N$, and Hodge index with $H^2=2N$ gives the bound.

If $D^2=0$, write $D=\sum n_iC_i$ with integral components. Nefness forces $C_i^2=C_i\cdot C_j=0$. Adjunction gives $p_a(C_i)=1$. An abelian variety contains no rational curve; hence the normalization of $C_i$ has genus one, its delta invariant is zero, and $C_i$ is smooth. After translation it is an elliptic subgroup. Different elliptic subgroup directions have positive intersection, since their addition map onto the surface is an isogeny. All components are therefore fibres of one elliptic quotient $q:A\to E$. We have $D=q^*D_0$ with $D_0$ effective, and connected fibres give $h^0(D)=h^0(D_0)=\deg D_0$. The restriction of $H$ to a fibre is globally generated and has positive degree. It cannot have degree one, whose sole section vanishes. Thus $H$ has degree at least two on each fibre, and $D\cdot H\ge2\deg D_0\ge2N$.
\end{proof}

We recall the precise theta input. The theta group $\mathcal G(L)$ is the central extension of $K(L)=\ker\lambda_L$ by $\mathbb G_m$ whose elements lift translations to $L$. Its commutator $e_L$ is a perfect alternating pairing. When $N$ is odd, $K(L)$ is \'etale of order $N^2$ and has complementary maximal isotropic subgroups of order $N$. Their theta extensions split. These facts, including descent of a linearized line, are proved in \citet[Theorem~8.10 and Corollary~8.11, pp.~116--117; Proposition~8.19, Corollary~8.20, Lemma~8.24 and Theorem~8.28, pp.~119--123]{EGM}.

\begin{lemma}[Theta weights and evaluation]\label{lem:evaluation}
Suppose $L$ is ample and globally generated, with odd $N=h^0(L)>1$. Set $E=F_{A*}L$ and $H=L^{(2)}$. Every nonzero representation of $\mathcal G(L)$ on which the centre acts by $z\mapsto z^r$, with $\gcd(r,N)=1$, has dimension divisible by $N$. The saturation $Q\subset E$ of the global evaluation image has rank at least three.
\end{lemma}
\begin{proof}
Diagonalize the action of one maximal isotropic level subgroup. The other subgroup permutes its $N$ character spaces through the characters $e_L(-,y)^r$. Coprimality makes this permutation transitive on all $N$ characters. All multiplicities are equal, proving divisibility. A representation of dimension $N$ and such a weight is irreducible by the same argument applied to a nonzero subrepresentation.

The theta action on $E$ has central weight one and covers translation by $F_A(x)$ for $x\in K(L)$. The image and saturation $Q$ are invariant. If their rank were $q=1$ or $2$, wedge evaluation would give a nonzero equivariant map
\[
 \bigwedge^q H^0(A,L)\longrightarrow H^0(B,\det Q).
\]
It is a global map to the reflexive determinant: the determinant sections defined off the finite singular set extend over it. Its image has weight $q$, so $h^0(\det Q)\ge N$. Lemma~\ref{lem:divisor} gives $c_1(Q)\cdot H\ge2N$. On the other hand $c_1(E)=2c_1(H)$ numerically and $H^2=2N$, so $\mu_H(E)=N$. Lemma~\ref{lem:stable} implies $c_1(Q)\cdot H<qN\le2N$, a contradiction.
\end{proof}

\begin{proposition}[Odd Cartier surjectivity]\label{prop:surjection}
If $A$ is nonordinary and $L$ satisfies the hypotheses of Lemma~\ref{lem:evaluation}, then Cartier multiplication induces a surjection
\[
 \beta:\bigwedge^2H^0(A,L)\longrightarrow H^0(B,H).
\]
\end{proposition}
\begin{proof}
Finite-flat duality for Frobenius, with an invariant volume form, gives a perfect pairing $E\otimes E\to H$ by Cartier trace of products. Locally this is the coefficient pairing in the truncated algebra with basis $1,x,y,xy$, so its perfection is also immediate from the Cartier coefficient formula. It is alternating on a nonordinary surface. Indeed Cartier of a square times the invariant top form is the square root factor times Cartier of that form, and Cartier kills the invariant top form precisely when the surface is nonordinary. Its Serre-dual map on $H^2(A,\OO_A)$ is the determinant of Frobenius on $H^1(A,\OO_A)$, which vanishes exactly when the latter is not bijective.

If its restriction to global sections vanished, the generic evaluation image would be isotropic in a four-dimensional symplectic space. Its dimension would be at most two, contrary to Lemma~\ref{lem:evaluation}. Thus $\beta$ is nonzero. It is theta equivariant, and $H^0(B,H)$ has dimension $N$ and central weight two. Lemma~\ref{lem:evaluation} makes this target irreducible; hence the image is the whole target.
\end{proof}

\begin{proposition}[Generic generation prime to six]\label{prop:rank-four}
Under the hypotheses of Lemma~\ref{lem:evaluation}, if $3\nmid N$, then evaluation $H^0(A,L)\otimes\OO_B\to E$ has generic rank four. Nonordinarity is not required.
\end{proposition}
\begin{proof}
Only rank three remains. Suppose $\rk Q=3$, put $T=E/Q$, and let $J=T^{**}$. This line has central weight one; $D=\det Q$ has weight three. Wedge evaluation and $\gcd(3,N)=1$ imply $h^0(D)\ge N$.

We first show that $J$ is ample. Pulling the quotient back and using the double dual gives a map
\[
 U=F_A^*E\otimes L^{-1}\longrightarrow F_A^*J\otimes L^{-1}
\]
that is surjective off finitely many points. Restrict to the normalization of any integral curve. The image is the target line minus an effective divisor. Since the restriction of $U$ is an extension of trivial bundles, every line quotient has nonnegative degree. Hence $F_A^*J\otimes L^{-1}$ has nonnegative degree on every curve and is nef. Tensoring with ample $L$ makes $F_A^*J$ ample. Finite surjective pullback detects ampleness, so $J$ is ample.

Now $h^0(J)>0$ and weight one implies $h^0(J)\ge N$. The divisor bound applied to $D,J$ gives
\[
 D\cdot H\ge2N,\qquad J\cdot H\ge2N,\qquad
 (D+J)\cdot H=4N.
\]
Both inequalities are equalities. Since $J$ is ample, $J^2=2h^0(J)\ge2N$, while Hodge index gives $J^2\le2N$. Its equality case yields $J\equiv H$. Numerical triviality gives a torsion N\'eron--Severi class by \citet[Theorem~6.3, pp.~52--54]{Kleiman}, and this group is torsion free for an abelian variety \citep[Corollaries~7.23 and 7.25, pp.~106--107]{EGM}. Therefore $J\otimes H^{-1}\in\Pic^0(B)$.

Tensoring by a homogeneous line leaves the theta commutator unchanged. Consequently, on $F_A(K(L))=K(H)$ the commutator of $J$ equals that of $H$, namely $e_L^2$. But the weight-one action on $J$ gives the commutator $e_L$. All values would satisfy $z=z^2$, hence be one. This contradicts nondegeneracy of $e_L$ and $N>1$.
\end{proof}

\section{Smooth curves in the complete system}\label{sec:smooth}

We prove existence of smooth members by a nodal deformation, so no separability form of Bertini is implicit in the argument.

\begin{lemma}[Odd globally generated systems have smooth members]\label{lem:smooth}
Every ample globally generated line $L$ on an abelian surface with odd $N=h^0(L)>1$ has a nonempty open of smooth connected members in $|L|$, each of genus $N+1$.
\end{lemma}
\begin{proof}
Choose a maximal isotropic subgroup $K_0\subset K(L)$ of order $N$. Splitting its theta extension gives a linearization and descent along the \'etale quotient $\eta:A\to A/K_0$. The descended ample line $M$ has $\chi(M)=\chi(L)/N=1$, so is principal. The precise descent statement is \citet[Theorem~8.10 and Corollary~8.11, pp.~116--117]{EGM}.

We give the theta-divisor alternative needed here. A principal theta divisor $\Theta$ has $\Theta^2=2$. If it has a single integral component, say $\Theta=nD$, adjunction makes $D^2$ an even integer and $n^2D^2=2$ forces $n=1$. The arithmetic genus is two. If $D$ were singular, its normalization would have genus at most one. Genus zero is impossible in an abelian variety. A map from an elliptic curve to an abelian variety, after translation, is a homomorphism and has smooth elliptic image; this also contradicts singularity and arithmetic genus two. Thus $D$ is smooth of genus two.

If $\Theta=\sum n_iD_i$ has more than one component, ampleness gives positive integers $\Theta\cdot D_i$ whose weighted sum is two. There are exactly two components, each of multiplicity one and each meeting $\Theta$ once. Nefness and even self-intersection imply
\[
 D_1^2=D_2^2=0,\qquad D_1\cdot D_2=1.
\]
The argument of Lemma~\ref{lem:divisor} makes both smooth elliptic translates. Intersection length one means one transverse intersection. Thus $\Theta$ is reduced nodal. Its \'etale inverse image is a smooth or reduced nodal member $C_0=(s_0=0)$ of $|L|$.

If there are nodes, global generation permits a section $s_1$ nonzero at all of them: avoid finitely many proper hyperplanes in $H^0(L)$. Consider the pencil $s_0+z s_1$. At a node its completed local equation is $xy+z u(x,y,z)=0$ with $u$ a unit. On the relative singular subscheme its two partial derivatives give $x,y\in(z)$. The equation then gives $z$ times a unit in the same ideal, so this singular subscheme has no component dominating the parameter trait. Away from the nodes the central fibre is smooth. Properness forces any putative generic singularity to specialize to the central fibre, which the preceding analysis excludes. The generic member is therefore smooth. If $C_0$ is smooth, the same conclusion follows from openness. Connectedness and the genus formula follow from Lemma~\ref{lem:cohomology}.
\end{proof}

\section{The two Heisenberg contraction matrices}\label{sec:matrices}

We separate the finite-dimensional algebra from its geometric realization. Let $\ell$ be any odd prime, choose a primitive $\ell$th root $\zeta\in k$, and let $W$ have basis $e_i$ for $i\in\mathbb Z/\ell$. The shift $S$ and clock $T$ act by
\[
 Se_i=e_{i+1},\qquad Te_i=\zeta^i e_i,
 \qquad TS=\zeta ST.
\]
Their finite Heisenberg group has centre $\mu_\ell$, acting on $W$ with weight one and on $W^{(2)}$ with weight two. These are the theta representations of a type $(1,\ell)$ polarization.

\begin{lemma}[Contraction formulas]\label{lem:matrices}
Every equivariant alternating map $\beta:\bigwedge^2W\to W^{(2)}$ has the form
\begin{equation}\label{eq:tensor}
 \beta(e_i,e_j)=c_{i-j}e_{(i+j)/2}^{(2)},\qquad
 c_0=0,\quad c_a=c_{-a}.
\end{equation}
Its dual contraction and multiplication matrix are
\begin{equation}\label{eq:matrices}
 B_y(i,j)=c_{i-j}y_{(i+j)/2},\qquad
 M_s(k,j)=c_{2(k-j)}s_{2k-j}.
\end{equation}
They satisfy $M_ss=0$ and
\begin{equation}\label{eq:incidence}
 B_ys=0\quad\Longleftrightarrow\quad M_s^t y=0.
\end{equation}
\end{lemma}
\begin{proof}
Clock equivariance forces the output index $h$ to satisfy $2h=i+j$. Shift equivariance makes the coefficient depend only on $i-j$. Alternation gives exactly the displayed conditions. Contracting the tensor with $y$ gives $B_y$. In the coefficient of output $e_k^{(2)}$ in $\beta(s,e_j)$, the section index is $2k-j$, giving $M_s$. The remaining identities are respectively $\beta(s,s)=0$ and the equality of the two ways of evaluating $y(\beta(s,v))$ for every $v$.
\end{proof}

\begin{proposition}[Universal maximal dual rank for six primes]\label{prop:certificate}
For $\ell\in\{3,5,7,11,13,17\}$, every nonzero coefficient vector in \eqref{eq:tensor} gives $\rk B_y=\ell-1$ for generic $y$. This is valid over every field of characteristic two.
\end{proposition}
\begin{proof}
Put $m=(\ell-1)/2$ and delete index zero from $B_y$. Its Pfaffian $P_0$ is a polynomial over $\FF_2[c_1,\ldots,c_m,y_0,\ldots,y_{\ell-1}]$. For an even subset $I$ of the remaining indices define $P(\varnothing)=1$, and, with $i$ its least element, use
\begin{equation}\label{eq:pfaff-recurrence}
 P(I)=\sum_{j\in I\setminus\{i\}}
 c_{i-j}y_{(i+j)/2}\,P(I\setminus\{i,j\}).
\end{equation}
This recurrence is the Pfaffian expansion; signs disappear in characteristic two. Thus it specifies the whole polynomial by finite additions and multiplications, with no specialization of the coefficients.

A support is a nonempty subset $D\subset\{1,\ldots,m\}$, indicating exactly which $c_a$ are nonzero. For every support in each listed prime, the recurrence produces a $y$-monomial whose coefficient becomes a single $c$-monomial when $c_a=0$ for $a\notin D$. Such a coefficient is nonzero at every point of that support torus. The exact enumeration has the following sizes:
\[
\begin{array}{c|r|r|r|r}
\ell&\text{formal matchings}&\text{surviving monomials}
 &\text{$y$-coefficients}&\text{supports}\\ \hline
3&1&1&1&1\\
5&3&3&3&3\\
7&15&15&12&7\\
11&945&811&273&31\\
13&10395&8427&1428&63\\
17&2027025&1465345&43263&255
\end{array}
\]
The accompanying program implements \eqref{eq:pfaff-recurrence} with memoized subsets and parity cancellation of monomial sets. It gives, for each support, the exponent vector of a witnessing $y$-monomial, the surviving $c$-monomial, and every monomial in its unrestricted coefficient. The recurrence and witness test are reproduced in Appendix~\ref{app:certificate}; they constitute an exhaustive finite polynomial verification, not field sampling.

The first cases can be checked by hand. At $\ell=3$, $P_0=c_1y_0$. At $\ell=5$,
\[
 P_0=c_1^2y_1y_4+c_2^2y_2y_3+c_1c_2y_0^2.
\]
The first or second term already works on every nonzero coefficient vector. At $\ell=7$, distinct $y$-coefficients include $c_1c_2^2$, $c_1^2c_3$, and $c_2c_3^2$. If all vanish, at most one $c_a$ is nonzero, and its pure-cube term survives. The same support test is fully enumerated for the larger primes. Consequently $P_0(c,y)$ is nonzero for each nonzero $c$. Its square is the corresponding determinant, so the alternating matrix has rank at least $\ell-1$. An odd-dimensional alternating matrix has rank at most $\ell-1$, proving equality.
\end{proof}

\section{Incompressibility and the radical map}\label{sec:dominance}

We prove that maximal dual rank forces a birational radical map. The only incompressibility statement needed is the known prime-degree case, for which Karpenko's diagonal-trace method has a short direct proof. Here $\SB(D)$ denotes the Severi--Brauer variety of right ideals of reduced dimension one in $D$; it becomes projective space over a splitting field. We write $\CH^j$ for codimension-$j$ cycles modulo rational equivalence.

\begin{lemma}[Prime-degree incompressibility]\label{lem:incompressible}
Let $D$ be a central division algebra of prime degree $\ell$ over a field $K$. Every rational self-map of $Y=\SB(D)$ is dominant, with no separability assumption.
\end{lemma}
\begin{proof}
Every closed point of $Y$ has degree divisible by $\ell$. Indeed, its residue field $E$ splits $D$, so the left $D$-module $E^\ell$ has $K$-dimension $\ell[E:K]$, divisible by $\dim_KD=\ell^2$. Thus every zero-cycle degree is divisible by $\ell$.

Set $n=\ell-1$ and let $\Gamma\in\CH^n(Y\times Y)$ be the closure of the graph of a rational self-map. Over an algebraic closure write
\[
 \overline\Gamma=\sum_{i=0}^n a_i h_1^i h_2^{n-i}.
\]
Here $h_1,h_2$ are the hyperplane classes on the two split factors. Its first projection is birational, so $a_0=1$. These split correspondences are mutually orthogonal idempotents under composition: intersection in the middle factor has degree one precisely when the indices agree. The coefficients of $\Gamma^{\circ(\ell-1)}$ are therefore $a_i^{\ell-1}$. Pulling back by the diagonal regular embedding gives a zero-cycle of degree
\[
 \sum_{i=0}^{\ell-1}a_i^{\ell-1}.
\]
Modulo $\ell$ this is the number of coefficients not divisible by $\ell$. It must be zero modulo $\ell$, and $a_0=1$ ensures that this number is positive. There are exactly $\ell$ coefficients, so all are nonzero modulo $\ell$. In particular $a_n\ne0$. But $a_n$ is the multiplicity of the second projection and would be zero for a nondominant map. This is the required contradiction. The argument uses ordinary correspondence composition and refined diagonal intersection, and works in arbitrary characteristic.
\end{proof}

\begin{lemma}[Dominance of an equivariant map]\label{lem:dominance}
Every Heisenberg-equivariant rational map
\[
 \kappa:\PP((W^{(2)})^*)\dashrightarrow\PP(W)
\]
is dominant.
\end{lemma}
\begin{proof}
The projective actions factor through $(\mathbb Z/\ell)^2$. Work over $K=k(a)((b))$ and twist both spaces by the torsor obtained by adjoining $A_0,B_0$ with $A_0^\ell=a$, $B_0^\ell=b$. The generators multiply $A_0$ and $B_0$ by $\zeta$, respectively. On matrix algebras combine these actions with $\operatorname{Ad}(S)$ and $\operatorname{Ad}(T)$. The invariant algebra contains
\[
 X_0=A_0T,\qquad Y_0=B_0^{-1}S,
 \qquad X_0^\ell=a,\quad Y_0^\ell=b^{-1},
 \quad X_0Y_0=\zeta Y_0X_0.
\]
Its $\ell^2$ monomials form a basis. With $u=Y_0^{-1}$ it is the cyclic algebra
\[
 D=\bigoplus_{i=0}^{\ell-1}K(a^{1/\ell})u^i,
 \qquad u^\ell=b,\quad u\lambda=\sigma(\lambda)u,
\]
where $\sigma(a^{1/\ell})=\zeta a^{1/\ell}$. This algebra is a division algebra: the extension $K(a^{1/\ell})/K$ is unramified of residue degree $\ell$ for the $b$-adic valuation, and the valuations of the summands $\lambda_i u^i$ have distinct fractional parts $i/\ell$. A nonzero sum has a unique lowest term. Products retain a nonzero lowest term, even when an exponent crosses $\ell$, so there are no zero divisors. Finite dimension then implies division.

Write $\alpha=[D]$. The projective cocycle of a central-weight-$r$ representation is the $r$th power of the weight-one cocycle. The two twisted projective spaces therefore have Brauer classes $-2\alpha$ and $\alpha$ and dimension $\ell-1$. Call them $Z$ and $Y=\SB(D)$. Equivariance descends $\kappa$ to $Z\dashrightarrow Y$. Over $K(Y)$ the class $\alpha$ vanishes, so $Z_{K(Y)}$ is projective space. This infinite field has a rational point in the nonempty open where $\kappa$ is defined. Such a point defines a rational map $Y\dashrightarrow Z$ whose composite with $\kappa$ is defined. Lemma~\ref{lem:incompressible} makes that composite dominant, and hence makes $\kappa$ dominant. Dominance is unchanged by twisting and extension of the ground field, so the original map is dominant.
\end{proof}

\begin{theorem}[A birational radical map]\label{thm:birational}
For any odd prime $\ell$, suppose the tensor \eqref{eq:tensor} has generic dual rank $\ell-1$. Its radical map
\[
 \kappa:\PP((W^{(2)})^*)\dashrightarrow\PP(W),
 \qquad [y]\longmapsto\ker B_y,
\]
is birational. The matrix $M_s$ has generic rank $\ell-1$. The maximal Pfaffians of $B_y$ have no nonconstant common factor; consequently the nonmaximal-rank locus has codimension at least two.
\end{theorem}
\begin{proof}
On the nonempty maximal-rank open $U$, the radical is a line and defines an equivariant map. Lemma~\ref{lem:dominance} makes it dominant. Its fibre over $[s]$, by \eqref{eq:incidence}, is the open subset
\[
 U\cap\PP(\ker M_s^t).
\]
The generic fibre is nonempty and zero dimensional, because source and target have equal dimension. An open subset of a projective linear space has that dimension only when the space is a point. Thus $\rk M_s=\ell-1$ generically. Kernel minors give a rational inverse. This proves birationality as a map of function fields, including separability, rather than merely injectivity on geometric points.

Let $g$ be the greatest common divisor of the maximal Pfaffians. Their ideal is Heisenberg invariant, so uniqueness of the greatest common divisor makes $kg$ a one-dimensional representation. Every such character is trivial on the commutator centre. That centre scales $y$ with weight $-2$, and hence scales a homogeneous polynomial by weight $-2\deg g$ under substitution. Therefore $\ell\mid2\deg g$. But the Pfaffians have degree $(\ell-1)/2$, so $0\le\deg g<\ell$ and $\deg g=0$. A codimension-one component in their common zero locus would yield a common factor; none exists.
\end{proof}

\section{Generic bijectivity and the cubic example}\label{sec:conclusion}

We now convert the radical calculation into the geometric stable rank, and work out the smallest system explicitly.

\begin{proposition}[The Euler determinant is nonzero]\label{prop:generic}
Under the hypotheses of Theorem~\ref{thm:conditional}, maximal generic rank of $B_y$ implies that $\Phi:F_X^*\Omega_X^1\to\Omega_X^1$ is generically invertible.
\end{proposition}
\begin{proof}
Let $P(y)=(P_i(y))$ be the vector of maximal Pfaffians. It lies in $\ker B_y$, and is nonzero on the maximal-rank open. The identity follows either from Pfaffian expansion or by pairing the matchings in each component of $B_yP(y)$; each resulting monomial occurs twice. Its coordinates have degree $(\ell-1)/2$. Theorem~\ref{thm:birational} identifies $[s]=[P(y)]$ birationally with $[y]$, with $\ker M_s^t=ky$ on suitable dense opens.

Consider the polynomial
\[
 \Delta(y)=\sum_i y_iP_i(y)^2.
\]
It has degree $\ell$. In characteristic two,
\[
 \frac{\partial\Delta}{\partial y_i}=P_i(y)^2.
\]
Since the vector $P$ is nonzero, $\Delta$ is not identically zero. Its nonvanishing open meets the common birational open. At every point of this intersection, \eqref{eq:criterion} makes $\Phi$ invertible. The determinant is a regular section of $\OO_X(\ell)$, since $\det\Omega_X^1=\OO_X(-\ell)$ and Frobenius doubles its degree. Thus its nonvanishing locus is nonempty and open. The polynomial $\Delta$ is used on the dual parameter space; no identification with the determinant polynomial on the original space is asserted in general.
\end{proof}

\begin{proof}[Proof of Theorems~\ref{thm:main} and \ref{thm:conditional}]
Proposition~\ref{prop:surjection} gives the nonzero theta-equivariant Cartier tensor. In a theta basis Lemma~\ref{lem:matrices} gives \eqref{eq:tensor}. Proposition~\ref{prop:certificate} supplies maximal generic dual rank for the six primes, while Theorem~\ref{thm:conditional} assumes that rank for an arbitrary odd prime. Proposition~\ref{prop:generic} gives a nonempty Euler-bijective open. Lemma~\ref{lem:smooth} gives a nonempty smooth-member open. They meet because $\PP(W)$ is integral. On their intersection $Q_C$ is Frobenius bijective of dimension $\ell-1$, and Lemma~\ref{lem:cohomology} yields all three formulas. Algebraic closure gives the stated interpretation over a perfect field.
\end{proof}

\begin{proposition}[The type $(1,3)$ system]\label{prop:cubic}
For $\ell=3$, the tensor $\beta$ is an isomorphism. In theta coordinates on $X=\PP^2$, the nonbijective locus of $\Phi$ is the smooth Fermat cubic
\[
 s_0^3+s_1^3+s_2^3=0.
\]
The determinant vanishes to order one there. On a dense open of this cubic, the first and stable ranks of $\Phi$ are both one.
\end{proposition}
\begin{proof}
There is one coefficient $c=c_1=c_2$, and it is nonzero by surjectivity. Thus $\beta(e_i,e_j)=c e_h^{(2)}$ when $\{i,j,h\}=\{0,1,2\}$, an isomorphism. On $s_0=1$, put $s=(1,t,u)$. Then
\[
 M_s=c\begin{pmatrix}0&u&t\\u&0&1\\t&1&0\end{pmatrix}.
\]
The Euler frame of $s^\perp$ is $(ta+ub,a,b)$, while its Frobenius pullback is $(t^2a+u^2b,a,b)$. Applying the transpose matrix gives
\[
 \Phi=c\begin{pmatrix}ut^2&1+u^3\\1+t^3&tu^2\end{pmatrix},
 \qquad \det\Phi=c^2(1+t^3+u^3).
\]
The other charts give the homogeneous cubic above. Its partial derivatives are $s_0^2,s_1^2,s_2^2$, with no common projective zero, so it is smooth and the determinant has multiplicity one.

On the cubic the matrix factors as
\[
 \Phi=c\binom{t^2}{u^2}(u\quad t).
\]
The scalar controlling the next semilinear iterate is a nonzero constant times
\[
 ut^4+tu^4=ut(t^3+u^3)=ut.
\]
For $tu\ne0$ the image line is Frobenius bijective, proving both rank assertions. This is a computation on the full universal Euler bundle, including coefficients defining singular divisors.
\end{proof}

\section{The remaining rank question}\label{sec:open}

The preceding results isolate a finite-dimensional obstruction for every odd polarization prime. The open question is whether, for every odd prime $\ell$, every nonzero Heisenberg-equivariant alternating tensor \eqref{eq:tensor} has generic dual rank $\ell-1$. A weaker question sufficient for geometry asks the same only for tensors realized by Cartier multiplication on the globally generated nonordinary abelian surfaces under consideration. Surjectivity of $\beta$ does not settle either question, and generic generation of $F_{A*}L$ when $3\nmid\ell$ does not settle it either.

A positive answer gives the geometric rank formula through Theorem~\ref{thm:conditional}, without any new birationality argument. The certificates here prove only the six listed primes. The gradient identity for $\Delta$ shows that its zero locus is smooth on the dual maximal-rank open, but does not show that it meets the open on which the radical map is an isomorphism: a component could be contracted. Thus common-wall, local-inertia, Lang-cover and full-monodromy questions require additional arguments. No such conclusion is part of this paper.

\appendix
\section{Exact finite certificate}\label{app:certificate}

This appendix specifies the complete finite verification behind Proposition~\ref{prop:certificate}. Monomials are represented by exponent vectors in the $\ell+m$ variables. Multiplication adds exponent vectors. Addition over $\FF_2$ toggles membership in a set, so repeated monomials cancel. Apply recurrence \eqref{eq:pfaff-recurrence} to every required even subset, storing each result once. Group the resulting monomials of $P_0$ by their $y$-exponent vector. For a fixed group write its coefficient as $\sum_{a\in T}c^a$. For each nonempty support $D$, retain those exponent vectors with $\operatorname{supp}(a)\subset D$. Accept this group as a witness precisely when one vector remains.

These instructions are a terminating exhaustive algorithm on the displayed finite inputs. Its output contains one witness for every support; each witness can be checked by the same recurrence restricted to the chosen $y$-coefficient. The implementation packs each exponent into four bits. Every exponent is bounded by the Pfaffian degree, at most eight, so addition of packed exponents causes no carry into another variable. This encoding is an optimization of exact polynomial arithmetic, not a modular reduction of the exponents. The number of nonempty supports is $2^m-1$, accounting for the last column of the table. Once a witness is found, its unique surviving monomial is a product of nonzero field elements, which proves nonvanishing on that entire support torus over any characteristic-two field.

\section*{Sources and method}
The classical inputs are credited at their exact locations in the text. The stability, divisor, smoothing, correspondence and Euler-map arguments are written in full, so the proof does not depend on a companion manuscript. The \emph{Prime Heisenberg Pfaffian certificate}, its output and build instructions are available through \url{\workbenchurl}; permanent records are at \url{\recordurl}. No general-prime conclusion is inferred from its finite scope.

The mathematics and text were produced by GPT-6 Astra Pro in ChatGPT Pro mode. A statement-led author attempt, with disclosed earlier exposure to related abelian-surface arguments, was frozen before the primitive proofs were read; it was partial. The symbolic Pfaffian computation was reconstructed before the earlier programs were read. The printed rank proofs and the cited primary proof passages were then checked. The earlier two Pfaffian programs were read and rerun, and their exact counts and recorded witnesses agreed with the separate reconstruction. A separate GPT-6 Astra Pro review of all proofs and an exact certificate rerun found no defect. A separate GPT-6 Astra Pro reconstruction from the statements alone proved both main theorems and all supporting assertions; its frozen proofs and two further exact certificate programs were compared in full. It used Karpenko's published incompressibility theorem explicitly. These are same-model checks, not human or formal verification. Reasoning-effort metadata is not exposed. No human has checked the mathematics. The new text and code are dedicated to the public domain under CC0 1.0.

\begin{thebibliography}{99}
\bibitem[Catanese and Sernesi(2022)]{CataneseSernesi}
Catanese, F., \& Sernesi, E. (2022, December 30). \emph{The Hesse pencil and polarizations of type $(1,3)$ on Abelian surfaces} (Version 1). arXiv:2212.14877. \url{https://arxiv.org/abs/2212.14877v1}.
\bibitem[Edixhoven et~al.(n.d.)]{EGM}
Edixhoven, B., van der Geer, G., \& Moonen, B. (n.d.). \emph{Abelian varieties: Preliminary version of the first chapters}. Author manuscript. \url{https://van-der-geer.nl/~gerard/AV.pdf}.
\bibitem[Grothendieck(n.d.)]{Grothendieck}
Grothendieck, A. (n.d.). \emph{Notes on crystalline cohomology}, cote~4, \S8.5.8, archive sheets~7--8. Institut Montpelli\'erain Alexander Grothendieck. \url{https://grothendieck.umontpellier.fr/4.pdf}.
\bibitem[Illusie(1979)]{Illusie}
Illusie, L. (1979). Complexe de de Rham--Witt et cohomologie cristalline. \emph{Annales scientifiques de l'\'Ecole normale sup\'erieure, 12}(4), 501--661. \url{https://doi.org/10.24033/asens.1374}.
\bibitem[Karpenko(2011)]{Karpenko}
Karpenko, N. A. (2011, July 9). \emph{Upper motives of algebraic groups and incompressibility of Severi--Brauer varieties}. Updated author preprint, originally dated April 18, 2009. \url{https://sites.ualberta.ca/~karpenko/publ/upperinn.pdf}.
\bibitem[Kleiman(2005)]{Kleiman}
Kleiman, S. L. (2005, April 1). \emph{The Picard scheme} (Version 1). arXiv:math/0504020. \url{https://arxiv.org/abs/math/0504020v1}.
\bibitem[Oda(1969)]{Oda}
Oda, T. (1969). The first de Rham cohomology group and Dieudonn\'e modules. \emph{Annales scientifiques de l'\'Ecole normale sup\'erieure, 2}(1), 63--135. \url{https://doi.org/10.24033/asens.1175}.
\bibitem[OpenAI(2026)]{OpenAI}
OpenAI. (2026, September 24). \emph{A cubic lower bound for border determinantal complexity of the permanent}. Research preprint. \url{https://github.com/openai/math}.
\bibitem[Sun(2006)]{Sun}
Sun, X. (2006, August 18). \emph{Stability of direct images under Frobenius morphism} (Version 2). arXiv:math/0608043. \url{https://arxiv.org/abs/math/0608043v2}.
\bibitem[Zou(2026)]{Zou}
Zou, H. (2026, September 29). \emph{Quasi-$F$-split primitive symplectic varieties in positive characteristic}. Author preprint. \url{https://www.math.uni-bielefeld.de/~hzou/papers/qfpsv/}.
\end{thebibliography}
\end{document}
